\documentclass[12pt]{article}
\usepackage[utf8]{inputenc}
\usepackage[spanish]{babel}
\usepackage{booktabs}
\usepackage{listings}
\usepackage{xcolor}
\usepackage{geometry}
\usepackage{amsmath}
\usepackage{array}
\usepackage{tabularx}
\usepackage{titlesec}
\usepackage{fancyhdr}
\usepackage{graphicx}
\geometry{margin=2.5cm}
\setlength{\headheight}{14pt}
\addtolength{\topmargin}{-2pt}

\pagestyle{fancy}
\fancyhf{}
\rhead{\small Compiladores --- UANL}
\lhead{\small Equipo: No Compila}
\rfoot{\thepage}

\lstset{
    basicstyle=\ttfamily\small,
    numberstyle=\tiny\color{gray},
    numbers=left,
    frame=single,
    backgroundcolor=\color{gray!8},
    rulecolor=\color{gray!40},
    breaklines=true,
    showstringspaces=false,
    keywordstyle=\color{blue!70!black}\bfseries,
    commentstyle=\color{green!50!black}\itshape,
    stringstyle=\color{red!60!black}
}

\begin{document}

% ══════════════════════════════════════════
%  PORTADA
% ══════════════════════════════════════════
\begin{titlepage}
\centering
\vspace*{1cm}

{\Large \textbf{Universidad Autónoma de Nuevo León}} \\[0.4cm]
{\large Facultad de Ciencias Físico Matemáticas} \\[2cm]

{\Huge \textbf{MiniPeka}} \\[0.3cm]
{\Large Compilador de un lenguaje tipo calculadora con funciones} \\[2cm]

{\large \textbf{Proyecto Integrador Académico (PIA)}} \\[0.2cm]
{\large Compiladores} \\[0.2cm]
{\large \textbf{Entregable 1}} \\[2cm]

\begin{tabular}{ll}
\textbf{Equipo:}      & No Compila \\[0.3cm]
\textbf{Integrantes:} & Regina Hernández Rodríguez \\
                      & César Oziel Guajardo Rodríguez \\
                      & Johan Joel Rodríguez Medina \\
                      & Rodrigo López Escobedo \\
                      & Diego Ángel Torres Azua \\
                      & Brandon Jahir Reyna Marroquín \\[0.5cm]
\textbf{Profesor:}    & Luis Ángel Gutiérrez Rodríguez \\[0.3cm]
\textbf{Fecha:}       & Mayo 2026 \\
\end{tabular}

\vfill
\end{titlepage}

% ══════════════════════════════════════════
%  TABLA DE CONTENIDO
% ══════════════════════════════════════════
\tableofcontents
\newpage

% ══════════════════════════════════════════
%  1. DEFINICIÓN DEL LENGUAJE
% ══════════════════════════════════════════
\section{Definición del lenguaje: MiniPeka}

MiniPeka es un lenguaje de programación simple diseñado con fines educativos
para ilustrar el proceso completo de compilación. Su sintaxis está inspirada
en lenguajes imperativos como C y Python, simplificada para facilitar la
implementación de cada fase del compilador. MiniPeka soporta declaración de
variables, operaciones aritméticas, estructuras de control básicas, definición
de funciones y entrada/salida por consola.

\subsection{Características principales}

\begin{itemize}
    \item \textbf{Tipado estático:} todas las variables deben declararse con
          su tipo antes de usarse. Los tipos disponibles son \texttt{int},
          \texttt{float} y \texttt{bool}.

    \item \textbf{Funciones con parámetros:} el lenguaje permite definir
          funciones con parámetros tipados y valor de retorno explícito.

    \item \textbf{Estructuras de control:} soporta \texttt{if}/\texttt{else}
          y el ciclo \texttt{while}.

    \item \textbf{Operadores aritméticos y relacionales:} suma, resta,
          multiplicación, división, y comparadores como \texttt{==},
          \texttt{!=}, \texttt{<}, \texttt{>}.

    \item \textbf{Instrucción de salida:} la función \texttt{print()}
          permite mostrar valores en consola.

    \item \textbf{Comentarios:} se soportan comentarios de una línea con
          el símbolo \texttt{\#}.
\end{itemize}

\subsection{Ejemplo representativo del lenguaje}

El siguiente fragmento muestra las capacidades principales de MiniPeka:

\begin{lstlisting}[caption={Ejemplo general de MiniPeka}]
# Declaracion de variables
int x = 5;
int y = 3;
float resultado = 0.0;

# Funcion que suma dos enteros
fun suma(int a, int b) : int {
    return a + b;
}

# Uso de la funcion y estructura de control
int total = suma(x, y);

if (total > 6) {
    print(total);
} else {
    print(0);
}

# Ciclo while
int i = 0;
while (i < total) {
    i = i + 1;
}
\end{lstlisting}

% ══════════════════════════════════════════
%  2. GRAMÁTICA PRELIMINAR
% ══════════════════════════════════════════
\section{Gramática preliminar}

La gramática de MiniPeka está definida como una gramática libre de contexto
(CFG). Se utiliza la notación BNF extendida donde \texttt{|} indica
alternativas, los elementos entre corchetes \texttt{[ ]} son opcionales y
los elementos entre llaves \texttt{\{ \}} pueden repetirse cero o más veces.

\subsection{Producciones principales}

\shorthandoff{<>}
\begin{align*}
\langle program \rangle      &\to \{ \langle statement \rangle \} \\[4pt]
\langle statement \rangle    &\to \langle decl \rangle
                              \mid \langle assign \rangle
                              \mid \langle if\_stmt \rangle
                              \mid \langle while\_stmt \rangle
                              \mid \langle fun\_def \rangle
                              \mid \langle print\_stmt \rangle
                              \mid \langle return\_stmt \rangle \\[4pt]
\langle decl \rangle         &\to \langle type \rangle\ \textbf{id}
                                   \ \texttt{=}\ \langle expr \rangle
                                   \ \texttt{;} \\[4pt]
\langle assign \rangle       &\to \textbf{id}\ \texttt{=}
                                   \ \langle expr \rangle\ \texttt{;} \\[4pt]
\langle if\_stmt \rangle     &\to \texttt{if}\ \texttt{(}
                                   \langle expr \rangle \texttt{)}
                                   \ \langle block \rangle
                                   \ [ \texttt{else}\ \langle block \rangle ] \\[4pt]
\langle while\_stmt \rangle  &\to \texttt{while}\ \texttt{(}
                                   \langle expr \rangle \texttt{)}
                                   \ \langle block \rangle \\[4pt]
\langle fun\_def \rangle     &\to \texttt{fun}\ \textbf{id}\ \texttt{(}
                                   \langle params \rangle \texttt{)}
                                   \ \texttt{:}\ \langle type \rangle
                                   \ \langle block \rangle \\[4pt]
\langle block \rangle        &\to \texttt{\{}
                                   \ \{ \langle statement \rangle \}
                                   \ \texttt{\}} \\[4pt]
\langle params \rangle       &\to \langle type \rangle\ \textbf{id}
                                   \ \{ \texttt{,}\ \langle type \rangle
                                   \ \textbf{id} \}
                                   \mid \varepsilon \\[4pt]
\langle print\_stmt \rangle  &\to \texttt{print}\ \texttt{(}
                                   \langle expr \rangle \texttt{)}
                                   \ \texttt{;} \\[4pt]
\langle return\_stmt \rangle &\to \texttt{return}\ \langle expr \rangle
                                   \ \texttt{;} \\[4pt]
\langle expr \rangle         &\to \langle expr \rangle\ \langle op \rangle
                                   \ \langle term \rangle
                                   \mid \langle term \rangle \\[4pt]
\langle term \rangle         &\to \textbf{id}
                                   \mid \textbf{num}
                                   \mid \texttt{(}\ \langle expr \rangle
                                   \ \texttt{)}
                                   \mid \langle fun\_call \rangle \\[4pt]
\langle fun\_call \rangle    &\to \textbf{id}\ \texttt{(}
                                   \langle args \rangle \texttt{)} \\[4pt]
\langle args \rangle         &\to \langle expr \rangle
                                   \ \{ \texttt{,}\ \langle expr \rangle \}
                                   \mid \varepsilon \\[4pt]
\langle op \rangle           &\to \texttt{+} \mid \texttt{-}
                                   \mid \texttt{*} \mid \texttt{/}
                                   \mid \texttt{==} \mid \texttt{!=}
                                   \mid \texttt{<} \mid \texttt{>} \\[4pt]
\langle type \rangle         &\to \texttt{int} \mid \texttt{float}
                                   \mid \texttt{bool}
\end{align*}
\shorthandon{<>}

% ══════════════════════════════════════════
%  3. IDENTIFICACIÓN DE TOKENS
% ══════════════════════════════════════════
\section{Identificación de tokens}

Los tokens son las unidades léxicas mínimas del lenguaje. El analizador
léxico de MiniPeka reconoce las siguientes categorías, cada una definida
mediante una expresión regular. Las palabras reservadas tienen prioridad
sobre el token \texttt{ID}: si una cadena coincide con ambos patrones,
el lexer la clasifica como palabra reservada. Los comentarios son
reconocidos y descartados por el scanner sin generar ningún token.

\begin{table}[h!]
\centering
\renewcommand{\arraystretch}{1.25}
\begin{tabular}{@{}p{2.8cm} p{4.5cm} p{5.0cm}@{}}
\toprule
\textbf{Token} & \textbf{Expresión regular} & \textbf{Ejemplos} \\
\midrule
\texttt{INT}        & \texttt{int}
                    & \texttt{int} \\
\texttt{FLOAT\_T}   & \texttt{float}
                    & \texttt{float} \\
\texttt{BOOL}       & \texttt{bool}
                    & \texttt{bool} \\
\texttt{IF}         & \texttt{if}
                    & \texttt{if} \\
\texttt{ELSE}       & \texttt{else}
                    & \texttt{else} \\
\texttt{WHILE}      & \texttt{while}
                    & \texttt{while} \\
\texttt{FUN}        & \texttt{fun}
                    & \texttt{fun} \\
\texttt{RETURN}     & \texttt{return}
                    & \texttt{return} \\
\texttt{PRINT}      & \texttt{print}
                    & \texttt{print} \\
\texttt{ID}         & \texttt{[a-zA-Z\_][a-zA-Z0-9\_]*}
                    & \texttt{x}, \texttt{suma}, \texttt{total} \\
\texttt{NUM\_INT}   & \texttt{[0-9]+}
                    & \texttt{5}, \texttt{42}, \texttt{100} \\
\texttt{NUM\_FLOAT} & \texttt{[0-9]+\textbackslash.[0-9]+}
                    & \texttt{3.14}, \texttt{0.5} \\
\texttt{ASSIGN}     & \texttt{=}
                    & \texttt{=} \\
\texttt{PLUS}       & \texttt{+}
                    & \texttt{+} \\
\texttt{MINUS}      & \texttt{-}
                    & \texttt{-} \\
\texttt{TIMES}      & \texttt{*}
                    & \texttt{*} \\
\texttt{DIVIDE}     & \texttt{/}
                    & \texttt{/} \\
\texttt{EQ}         & \texttt{==}
                    & \texttt{==} \\
\texttt{NEQ}        & \texttt{!=}
                    & \texttt{!=} \\
\texttt{LT}         & \texttt{<}
                    & \texttt{<} \\
\texttt{GT}         & \texttt{>}
                    & \texttt{>} \\
\texttt{LPAREN}     & \texttt{(}
                    & \texttt{(} \\
\texttt{RPAREN}     & \texttt{)}
                    & \texttt{)} \\
\texttt{LBRACE}     & \texttt{\{}
                    & \texttt{\{} \\
\texttt{RBRACE}     & \texttt{\}}
                    & \texttt{\}} \\
\texttt{SEMI}       & \texttt{;}
                    & \texttt{;} \\
\texttt{COMMA}      & \texttt{,}
                    & \texttt{,} \\
\texttt{COLON}      & \texttt{:}
                    & \texttt{:} \\
\texttt{COMMENT}    & \texttt{\#.*}
                    & \texttt{\# esto es un comentario} \\
\bottomrule
\end{tabular}
\caption{Tabla de tokens del lenguaje MiniPeka}
\end{table}

% ══════════════════════════════════════════
%  4. EJEMPLOS DE ENTRADA
% ══════════════════════════════════════════
\section{Ejemplos de entrada válidos e inválidos}

\subsection{Entradas válidas}

Los siguientes fragmentos son sintáctica y semánticamente correctos en
MiniPeka y deben ser procesados sin errores por el compilador.

\begin{lstlisting}[caption={Ejemplo válido 1: operaciones aritméticas básicas}]
int a = 10;
int b = 4;
int c = a + b * 2;
print(c);
\end{lstlisting}

\begin{lstlisting}[caption={Ejemplo válido 2: función con retorno y condicional}]
fun maximo(int x, int y) : int {
    if (x > y) {
        return x;
    } else {
        return y;
    }
}

int resultado = maximo(8, 3);
print(resultado);
\end{lstlisting}

\begin{lstlisting}[caption={Ejemplo válido 3: ciclo while con acumulador}]
int contador = 0;
int suma = 0;

while (contador < 5) {
    suma = suma + contador;
    contador = contador + 1;
}

print(suma);
\end{lstlisting}

\subsection{Entradas inválidas}

Los siguientes fragmentos contienen errores que el compilador debe
detectar y reportar con un mensaje claro.

\begin{lstlisting}[caption={Error léxico: carácter inválido}]
int x = 5@;   # '@' no pertenece al alfabeto de MiniPeka
\end{lstlisting}

\begin{lstlisting}[caption={Error sintáctico: paréntesis no balanceado}]
if (x > 3 {   # falta el ')' de cierre
    print(x);
}
\end{lstlisting}

\begin{lstlisting}[caption={Error sintáctico: instrucción sin punto y coma}]
int z = 10    # falta ';' al final de la declaracion
print(z);
\end{lstlisting}

\begin{lstlisting}[caption={Error semántico: variable no declarada}]
int z = w + 1;   # 'w' no ha sido declarada
\end{lstlisting}

\begin{lstlisting}[caption={Error semántico: tipos incompatibles}]
int x = 3.14;   # no se puede asignar float a int sin conversion
\end{lstlisting}

% ══════════════════════════════════════════
%  5. ALCANCE DEL PROYECTO
% ══════════════════════════════════════════
\section{Alcance del proyecto}

\subsection{Lo que el proyecto SÍ incluye}

\begin{itemize}
    \item Analizador léxico completo con manejo de errores léxicos y
          reporte de línea y columna del error.

    \item Analizador sintáctico basado en gramática LL con manejo de
          errores sintácticos y mensajes descriptivos.

    \item Generación del árbol de sintaxis abstracta (AST) para cualquier
          programa válido en MiniPeka.

    \item Tabla de símbolos con soporte de scopes anidados (global y local
          por función).

    \item Análisis semántico: verificación de tipos, variables declaradas
          antes de su uso y compatibilidad de tipos en asignaciones.

    \item Generación de código de tres direcciones (TAC) para expresiones
          aritméticas, asignaciones y estructuras de control.

    \item Al menos una optimización básica: propagación de constantes.

    \item Pruebas de funcionamiento con al menos cinco ejemplos válidos
          y cinco inválidos.

    \item Documentación técnica completa de cada fase del compilador.

    \item Video explicativo de 5 minutos demostrando el funcionamiento.
\end{itemize}

\subsection{Lo que el proyecto NO incluye}

\begin{itemize}
    \item Generación de código ejecutable (código máquina o ensamblador).

    \item Tipos de datos compuestos como arreglos, structs o cadenas de
          texto.

    \item Recursión directa o indirecta entre funciones.

    \item Forma SSA (considerada opcional según el PIA).

    \item Manejo de excepciones o errores en tiempo de ejecución.

    \item Interfaz gráfica de usuario.

    \item Operadores lógicos (\texttt{and}, \texttt{or}, \texttt{not}).
\end{itemize}

% ══════════════════════════════════════════
%  6. CRONOGRAMA Y ROLES
% ══════════════════════════════════════════
\section{Cronograma y roles del equipo}

\subsection{Roles del equipo}

\begin{table}[h!]
\centering
\renewcommand{\arraystretch}{1.25}
\begin{tabular}{@{}p{5.5cm} p{3.2cm} p{4.3cm}@{}}
\toprule
\textbf{Integrante} & \textbf{Rol principal} & \textbf{Responsabilidad} \\
\midrule
Regina Hernández Rodríguez
    & Líder de proyecto
    & Coordinación general, documentación y presentación final \\
César Oziel Guajardo Rodríguez
    & Desarrollador léxico
    & Implementación del scanner y manejo de errores léxicos \\
Johan Joel Rodríguez Medina
    & Desarrollador sintáctico
    & Implementación del parser y generación del AST \\
Rodrigo López Escobedo
    & Desarrollador semántico
    & Tabla de símbolos, chequeo de tipos y errores semánticos \\
Diego Ángel Torres Azua
    & Desarrollador de backend
    & Generación de TAC y optimización de código \\
Brandon Jahir Reyna Marroquín
    & Integración y pruebas
    & Integración de módulos, pruebas y validación del sistema \\
\bottomrule
\end{tabular}
\caption{Roles del equipo No Compila}
\end{table}

\subsection{Cronograma de actividades}

\begin{table}[h!]
\centering
\renewcommand{\arraystretch}{1.25}
\begin{tabular}{@{}p{1.5cm} p{4.5cm} p{4.2cm} p{2.3cm}@{}}
\toprule
\textbf{Semana} & \textbf{Actividad} &
\textbf{Responsable(s)} & \textbf{Entregable} \\
\midrule
1
    & Definición del lenguaje, gramática preliminar, tokens, ejemplos
      de entrada y alcance del proyecto
    & Todo el equipo
    & Entregable 1 \\
2
    & Implementación del scanner, expresiones regulares, manejo de
      errores léxicos y pruebas de tokenización
    & César, Brandon
    & Entregable 1 \\
3
    & Gramática formal completa, implementación del parser,
      árbol de derivación y generación del AST
    & Johan, Rodrigo
    & Entregable 2 \\
4
    & Tabla de símbolos, análisis semántico, verificación
      de tipos y anotación del AST
    & Rodrigo, Diego
    & Entregable 2 \\
5
    & Generación de TAC, optimización básica, integración
      de módulos y pruebas finales del sistema
    & Diego, Brandon, Regina
    & Entregable 2 \\
Cierre
    & Video explicativo, exposición del proyecto y
      documentación final
    & Todo el equipo
    & Presentación final \\
\bottomrule
\end{tabular}
\caption{Cronograma del proyecto MiniPeka}
\end{table}

% ══════════════════════════════════════════
%  7. DECLARACIÓN DE INTEGRIDAD ACADÉMICA
% ══════════════════════════════════════════
\section{Declaración de integridad académica}

Los integrantes del equipo \textbf{No Compila} declaramos que el presente
proyecto es de nuestra autoría original. Todo el código, análisis,
diagramas y documentación han sido desarrollados por los integrantes del
equipo. El uso de fuentes externas, bibliografía o herramientas de apoyo
será debidamente referenciado en el documento final. Nos comprometemos a
respetar el reglamento institucional de la Universidad Autónoma de Nuevo
León en materia de integridad académica y a no incurrir en plagio ni en
el uso indebido de código de terceros.

\vspace{1.5cm}

\noindent
\begin{tabular}{p{7cm} p{6cm}}
\rule{6cm}{0.4pt} & \rule{6cm}{0.4pt} \\
Regina Hernández Rodríguez & César Oziel Guajardo Rodríguez \\[1.5cm]
\rule{6cm}{0.4pt} & \rule{6cm}{0.4pt} \\
Johan Joel Rodríguez Medina & Rodrigo López Escobedo \\[1.5cm]
\rule{6cm}{0.4pt} & \rule{6cm}{0.4pt} \\
Diego Ángel Torres Azua & Brandon Jahir Reyna Marroquín \\
\end{tabular}

\vspace{1.5cm}
\noindent
Fecha: \underline{\hspace{5cm}} \hfill
Firma del profesor: \underline{\hspace{5cm}}

\end{document}